% =====================================================================
%  4. Geometric interpretation (overlay animation)
% =====================================================================
\section{Watching it move}

\begin{frame}{Gradient descent, one step at a time}
  \begin{columns}[T]
    \begin{column}{0.62\textwidth}
      \centering
      \only<1>{\includegraphics[height=0.70\textheight]{fig_gd_frame_0}}%
      \only<2>{\includegraphics[height=0.70\textheight]{fig_gd_frame_1}}%
      \only<3>{\includegraphics[height=0.70\textheight]{fig_gd_frame_2}}%
      \only<4>{\includegraphics[height=0.70\textheight]{fig_gd_frame_3}}%
      \only<5>{\includegraphics[height=0.70\textheight]{fig_gd_frame_4}}%
      \only<6>{\includegraphics[height=0.70\textheight]{fig_gd_frame_5}}%
      \only<7>{\includegraphics[height=0.70\textheight]{fig_gd_frame_6}}%
      \only<8>{\includegraphics[height=0.70\textheight]{fig_gd_frame_7}}%
    \end{column}
    \begin{column}{0.36\textwidth}
      \footnotesize
      The M3 quadratic, $\bA = \bigl(\begin{smallmatrix}3&1\\1&3\end{smallmatrix}\bigr)$, $\vb = (4,4)$,
      from $\x_0 = (-1.5,2)$, $\alpha = 0.2$.
      \begin{itemize}\footnotesize
        \item<1-> The orange arrow, $-\nabla f(\x_k)$, leaves at a right angle to the contour.
        \item<2-> The first step is long: $\norm{\nabla f}$ is big out here.
        \item<4-> The steps shrink on their own as $\nabla f\to\mathbf 0$, with $\alpha$ fixed.
        \item<6-> The path bends: on ellipses, $-\nabla f$ misses $\x^\star$.
        \item<8-> After 7 steps $f = -3.995$; the optimum is $-4$.
      \end{itemize}
    \end{column}
  \end{columns}
\end{frame}

\begin{frame}{The same descent on the line fit}
  \centering
  \only<1>{\includegraphics[width=\textwidth, height=0.6\textheight, keepaspectratio]{fig_ssr_frame_0}}%
  \only<2>{\includegraphics[width=\textwidth, height=0.6\textheight, keepaspectratio]{fig_ssr_frame_1}}%
  \only<3>{\includegraphics[width=\textwidth, height=0.6\textheight, keepaspectratio]{fig_ssr_frame_2}}%
  \only<4>{\includegraphics[width=\textwidth, height=0.6\textheight, keepaspectratio]{fig_ssr_frame_3}}%
  \only<5>{\includegraphics[width=\textwidth, height=0.6\textheight, keepaspectratio]{fig_ssr_frame_4}}%
  \only<6>{\includegraphics[width=\textwidth, height=0.6\textheight, keepaspectratio]{fig_ssr_frame_5}}%
  \only<7>{\includegraphics[width=\textwidth, height=0.6\textheight, keepaspectratio]{fig_ssr_frame_6}}%
  \only<8>{\includegraphics[width=\textwidth, height=0.6\textheight, keepaspectratio]{fig_ssr_frame_7}}%

  \raggedright\footnotesize
  The weight--height data of Section 1, both parameters free, $\alpha = 0.01$ from $(b, m) = (0, 1)$.
  Left: the point $(b, m)$ walks downhill across the SSR contours. Right: the line that point draws.
  Click through steps $0, 1, 2, 4, 15, 60, 250, 808$.
\end{frame}

\begin{frame}{The same walk on the SSR surface}
  \begin{columns}[T]
    \begin{column}{0.56\textwidth}
      \centering
      \includegraphics[width=\linewidth, height=0.76\textheight, keepaspectratio]{fig_ssr_surface}
    \end{column}
    \begin{column}{0.42\textwidth}
      \small
      $\SSR(b, m)$ is a bowl over the $(b, m)$ plane. Every point of the bowl is one line through the data.
      \begin{itemize}\footnotesize
        \item GD starts at $(0, 1)$ and rolls downhill. The dashed curve is its shadow on the contour map,
              the path of the previous slide.
        \item It drops into the valley in a few steps, then creeps along the valley floor, because the valley
              is long and narrow ($\kappa \approx 29$, see Section 8).
        \item After $808$ steps it stops at $(0.9486,\ 0.6411)$, the least-squares line $(0.9487,\ 0.6410)$.
      \end{itemize}
    \end{column}
  \end{columns}
\end{frame}

\begin{frame}{Round bowls are easy, long valleys are not}
  \centering
  % Round vs elongated contours: -grad points at the minimiser only for round ones
\begin{tikzpicture}[scale=0.8, >={Stealth[length=6pt]}]
  % round bowl
  \begin{scope}
    \foreach \r/\c in {1.6/6, 1.2/12, 0.8/20, 0.4/32}
      \fill[mBlue!\c] (0,0) circle (\r);
    \foreach \r in {0.4,0.8,1.2,1.6} \draw[mBlue!55, line width=0.4pt] (0,0) circle (\r);
    \fill[mGreen, draw=white] (0,0) circle (2.4pt);
    \draw[->, line width=1.6pt, mOrange] (1.13,1.13) -- (0.32,0.32);
    \fill[mOrange, draw=white] (1.13,1.13) circle (2.4pt);
    \node[font=\scriptsize, mInk, align=center] at (0,-2.05) {$\kappa = 1$: $-\nabla f$ aims at $\x^\star$};
  \end{scope}
  % elongated bowl  (semi-axes 2s and 0.6s)
  \begin{scope}[xshift=5.3cm]
    \foreach \s/\c in {1.6/6, 1.2/12, 0.8/20, 0.4/32}
      \fill[mBlue!\c] (0,0) ellipse ({2*\s} and {0.6*\s});
    \foreach \s in {0.4,0.8,1.2,1.6} \draw[mBlue!55, line width=0.4pt] (0,0) ellipse ({2*\s} and {0.6*\s});
    \fill[mGreen, draw=white] (0,0) circle (2.4pt);
    % point on the outer ellipse at parameter 50 degrees: (2.057, 0.735)
    \draw[mGrey, densely dotted, line width=0.7pt] (2.057,0.735) -- (0,0);
    \draw[->, line width=1.6pt, mOrange] (2.057,0.735) -- ++(-0.22,-0.88);
    \fill[mOrange, draw=white] (2.057,0.735) circle (2.4pt);
    \node[font=\scriptsize, mInk, align=center] at (0,-2.05) {$\kappa \gg 1$: $-\nabla f$ misses $\x^\star$};
  \end{scope}
\end{tikzpicture}


  \vspace{4pt}
  \raggedright
  \begin{itemize}\small
    \item \textbf{Round contours} ($\bA\propto\bI$): $-\nabla f$ points straight at $\x^\star$; one well-chosen step lands on it.
    \item \textbf{Stretched contours}: $-\nabla f$ is still perpendicular to the contour, so it \emph{misses} $\x^\star$.
          The iterates bounce across the valley and crawl along it.
    \item How stretched is measured by the \alert{condition number} $\kappa = \lmax/\lmin$. It returns in the step limit (Section~5),
          every rate (Section~7) and the zig-zag (Section~8).
  \end{itemize}
\end{frame}
